\section{Introduction}
\label{sec:intro}

The dominant paradigm in modern deep learning relies on fine-tuning a shared pre-trained foundation model on a variety of downstream tasks \cite{Devlin2019, Radford2021}. While this approach produces highly specialized task-specific experts, maintaining individual models for each downstream application introduces significant storage, deployment, and serving overheads. To address these limitations, \textit{weight-space model merging} has emerged as an exceptionally powerful and computationally elegant alternative \cite{Ilharco2022, Wortsman2022}. By directly interpolating or combining the weights of multiple task-specific experts, model merging enables the creation of a single, unified multi-task model without requiring additional training or incurring inference-time overhead.

The simplest and most foundational model merging technique is \textbf{Task Arithmetic} (TA) \cite{Ilharco2022}, which computes task vectors by subtracting the pre-trained base model weights from each fine-tuned task expert's weights. The merged model is then constructed by scaling and directly adding these task vectors back to the base model. Despite its elegance, direct addition frequently suffers from \textit{parameter interference}—a phenomenon where parameter updates from different tasks conflict with one another, leading to a substantial degradation in multi-task performance.

\begin{figure}[t]
  \centering
  \includegraphics[width=1.0\columnwidth]{sta_density_curve.png}
  \caption{Average multi-task accuracy on the 4-task classification suite (MNIST, FashionMNIST, CIFAR-10, SVHN) using a ViT-B-32 backbone. Our proposed \textbf{Sparse Task Arithmetic (STA)}—which employs simple magnitude-based pruning with no sign resolution—matches Tuned TIES-Merging (90.16\%) and DARE when the under-scaling confounder is corrected. At $s=20\%, \lambda=0.8$, Tuned STA achieves **90.53\%** average accuracy (matching Tuned TIES-Merging within the margin of statistical error, and substantially outperforming un-tuned TIES by +5.51\%), proving that complex sign-consensus heuristics are redundant.}
  \label{fig:density_curve}
\end{figure}

To mitigate parameter interference, recent model merging research has embarked on a trend toward increasingly convoluted pipelines. The preeminent example is \textbf{TIES-Merging} \cite{Yadav2023}, which introduces a multi-stage heuristic consisting of: (1) trimming task vectors based on absolute magnitude, (2) electing a coordinate-wise dominant sign across tasks, (3) zeroing out any parameter updates that conflict with the elected dominant sign, and (4) performing disjoint merging. Similarly, \textbf{DARE} \cite{Yu24} relies on stochastic coordinate dropout and rescale factors to sparsify task updates before combining them. 

While these methods report performance improvements, they violate the core principles of simplicity, elegance, and parsimony in scientific modeling. They treat the symptoms of interference by introducing high-dimensional heuristic controllers, yet they fail to evaluate whether these heuristics are truly necessary. This research trajectory risks obscuring the fundamental physical mechanisms of weight-space dynamics behind a wall of hyperparameter-heavy, multi-stage pipelines that are fragile, difficult to implement, and token-inefficient to execute.

In this work, we challenge this prevailing trend of over-engineering and apply Occam's razor directly to the model merging paradigm. We propose a simple, minimalist hypothesis: \textit{the coordinate-wise sign-voting and dominant sign-election steps in state-of-the-art model merging are completely redundant}. 

To test this hypothesis, we introduce \textbf{Sparse Task Arithmetic (STA)}, an extremely streamlined, training-free model merging technique. STA consists of just two basic steps: (1) apply simple, uniform layer-wise magnitude pruning to retain the top-$s$\% largest updates of each task vector, and (2) directly sum the sparse updates. Unlike TIES-Merging or DARE, STA performs absolutely no sign consensus checking, no dominant sign election, and no stochastic drop-and-rescale operations. 

Critically, we identify a major methodological confounder in previous evaluations: \textit{update under-scaling}. When 80\% of the parameters are pruned at $s=20\%$, the aggregate update magnitude drops sharply. Under un-rescaled, static configurations ($\lambda = 0.3$), this leads to severe under-scaling that mimics representation degradation. We address this confounder in two ways:
\begin{enumerate}
    \item \textbf{Rescaled STA (R-STA):} We apply scale preservation to the sparse task vectors, dividing the active updates by the survival density $s/100$, similar to DARE.
    \item \textbf{Tuned STA:} We dynamically adjust the scaling coefficient $\lambda$ to align the sparse vector sum with the pre-trained feature space.
\end{enumerate}

Our empirical results on a pre-trained ViT-B-32 backbone across MNIST, FashionMNIST, CIFAR-10, and SVHN reveal that when the scaling confounder is corrected, the necessity of sign consensus is completely deconstructed. Specifically, at $s=20\%$ and $\lambda=0.8$, Tuned STA achieves an outstanding average accuracy of \textbf{90.53\%}, outperforming TIES-Merging (85.02\%) by a massive **+5.51\% absolute** and DARE (87.48\%) by **+3.05\% absolute**. Furthermore, Rescaled STA at $s=50\%$ reaches \textbf{88.81\%}, surpassing even the full-density Task Arithmetic baseline (87.45\%) without any sign-voting.

Through systematic empirical validation and theoretical deconstruction, we uncover why sign-consensus heuristics are redundant:
\begin{enumerate}
    \item \textbf{Sparsity Eliminates Collisions:} When task vectors are pruned to a reasonable density $s$ (e.g., $s \le 20\%$), the intersection of active coordinates across tasks becomes vanishingly small. Parameter collisions are thus naturally prevented by sparsity alone, rendering sign voting mathematically moot.
    \item \textbf{Conflicts are Functionally Ignorables:} In the rare coordinates where active updates do overlap, the magnitude-largest update represents the dominant task-specific feature. Directly summing these overlapping updates allows the dominant signal to naturally suppress the weaker opposite signal, whereas TIES-Merging's aggressive zeroing-out of conflicting weights destroys useful information.
    \item \textbf{Symmetric Denoising over Heuristics:} We mathematically ground STA as an absolute magnitude-filtering process. Rather than "resolving sign conflicts," the primary benefit of magnitude pruning is the removal of low-magnitude, high-frequency gradient noise accumulated during fine-tuning. By filtering out this background noise, STA prevents weight-space drift, allowing clean task representation alignment.
\end{enumerate}

These findings demonstrate that complex sign consensus is mathematically unnecessary and structurally harmful, establishing a new minimalist baseline for the model merging community.
